\section{0x33\ 数学\ 组合计数}
\label{sec:0x33}

\subsection{排列组合}

\Highlight{排列数}：$A_{n}^m$，从 $n$ 个有标号元素中，取出 $m$ 个元素排成一列的方案数（需要考虑取数顺序）。

\[
A_n^m = \frac{n!}{(n - m)!} = n^{\underline{m}}
\]
\Highlight{组合数}：$C_n^m$ 或 $\binom{n}{m}$，从 $n$ 个有标号元素中，取出 $m$ 个元素组成集合的方案数（不需要考虑取数顺序）。

\[
\binom{n}{m} = \begin{cases}
\frac{n^{\underline{m}}}{m!} & m \in \mathbb{Z}, m \geq 0 \\
0 & m \in \mathbb{Z}, m < 0
\end{cases}
\]

\subsubsection{组合数常见恒等式}

\Highlight{阶乘展开式}：$n, m \in \mathbb{Z}$，$n \geq m \geq 0$

\[
\binom{n}{m} = \frac{n!}{m!(n - m)!}
\]
\Highlight{归纳恒等式}：$m \in \mathbb{Z}$

\[
\binom{n}{m} = \binom{n - 1}{m} + \binom{n - 1}{m - 1}
\]
\Highlight{对称恒等式}：$n, m \in \mathbb{Z}$，$n \geq 0$

\[
\binom{n}{m} = \binom{n}{n - m}
\]
\Highlight{吸收恒等式}：$m \in \mathbb{Z}$

\[
\begin{aligned}
\binom{n}{m} & = \frac{n}{m}\binom{n - 1}{m - 1} \\
\binom{n}{m} & = \frac{n}{n - m}\binom{n - 1}{m} \\
\end{aligned}
\]
\Highlight{上指标反转}：$m \in \mathbb{Z}$

\[
\binom{n}{m} = (-1)^m \binom{m - n - 1}{m}
\]
\Highlight{三项式版恒等式}：$m, k \in \mathbb{Z}$

\[
\binom{n}{m}\binom{m}{k} = \binom{n}{k}\binom{n - k}{m - k}
\]
\Highlight{二项式定理}：$n \in \mathbb{N}$

\[
(x + y)^n = \sum_{i = 0}^n \binom{n}{i}x^{i}y^{n - i}
\]

\begin{quoteblock}
一些常见特例：

\begin{itemize}
\item $x = 1$ 且 $y = a$：

\[
\sum_{i = 0}^{n} \binom{n}{i}a^i = (1 + a)^n
\]

\item $x = 1$ 且 $y = -1$：

\[
\sum_{i = 0}^n (-1)^i \binom{n}{i} = [n = 0]
\]
\end{itemize}
\end{quoteblock}

\Highlight{广义二项式定理}：$r \in \mathbb{R}$

\[
(x + y)^r = \sum_{k = 0}^{\infty} \frac{r^{\underline{k}}}{k!} x^{k}y^{r - k}
\]
\Highlight{平行求和法}：$m \in \mathbb{Z}$

\[
\sum_{0 \leq k \leq m} \binom{n + k}{k} = \binom{n + m + 1}{m}
\]
\Highlight{上指标求和法}：$m \in \mathbb{Z}$，$n \geq 0$

\[
\sum_{0 \leq k \leq n}\binom{k}{m} = \binom{n + 1}{m + 1}
\]
\Highlight{范德蒙德卷积}：$m \in \mathbb{Z}$

\[
\sum_k \binom{s}{k}\binom{t}{m - k} = \binom{s + t}{m}
\]
\Highlight{范德蒙德卷积变式}：$m \in \mathbb{Z}$

\[
\sum_k \binom{k}{s}\binom{m - k}{t} = \binom{m + 1}{s + t + 1}
\]
\Highlight{斜对角线求和法}：$n \in \mathbb{Z}$

\[
\sum_{i = 0}^n \binom{n - i}{i} = \mathrm{Fib}_{n + 1}
\]

\begin{quoteblock}
等价于长度为 $n - 1$ 的没有相邻 $1$ 的二进制串个数。
\end{quoteblock}

\Highlight{组合数带权和}：$n \in \mathbb{Z}$

\[
\sum_{i = 0}^n i\binom{n}{i} = n2^{n - 1}
\]
\Highlight{组合数带权平方和}：$n \in \mathbb{Z}$

\[
\sum_{i = 0}^n i^2\binom{n}{i} = n(n + 1)2^{n - 2}
\]

\subsubsection{插板法}

\Highlight{问题 1}：将 $n$ 个完全相同的元素分成 $k$ 组，保证每组至少有一个元素。等价于 $\sum_{i = 1}^k x_i = n$ 且 $\forall x_i \geq 1$ 的整数解个数。

相当于将 $k - 1$ 块板，插入 $n - 1$ 个间隙中。

\[
\binom{n - 1}{k - 1}
\]
\Highlight{问题 2}：将 $n$ 个完全相同的元素分成 $k$ 组，允许组内元素为空。等价于 $\sum_{i = 1}^k x_i = n$ 且 $\forall x_i \geq 0$ 的整数解个数。

给 $k$ 组分别补一个元素，转化为问题 1。

\[
\binom{n + k - 1}{k - 1}
\]
\Highlight{问题 3}：给 $n$ 个完全相同的元素分成 $k$ 组，保证第 $i$ 组至少分到 $a_i$ 个元素。

给第 $i$ 组提前安排 $a_i$ 个元素，转化为问题 2。

\[
\binom{n - \sum a_i + k - 1}{k - 1}
\]

\subsubsection{Lucas 定理}

\Highlight{Lucas 定理}：对于质数 $p$

\[
\binom{n}{m} \equiv \binom{\lfloor n / p \rfloor}{\lfloor m / p \rfloor}\binom{n \bmod p}{m \bmod p} \pmod p
\]
相当于是将 $n, m$ 转化为 $p$ 进制，将每一位上对应的组合数 $\binom{n_i}{m_i}$ 相乘后对 $p$ 取模。

\Highlight{组合数 mod 2}：当 $p = 2$ 时，$\binom{n}{m} \equiv \binom{\lfloor n / 2 \rfloor}{\lfloor m / 2 \rfloor}\binom{n \bmod 2}{m \bmod 2} \pmod 2$。

故 $\binom{n}{m} \bmod 2 = 1$，当且仅当\textbf{在二进制下 $m$ 是 $n$ 的子集}。

\subsection{组合数常见求法}

\subsubsection{求法 1：递推}

\begin{quoteblock}
\textbf{适用情况：对模数不做要求，可以承受平方时间的预处理（$n \leq 5000$）。}

时间复杂度：预处理 $\mathcal{O}(n^2)$，查询 $\mathcal{O}(1)$。
\end{quoteblock}

初值 $\binom{n}{0} = \binom{n}{n} = 1$，递推式 $\binom{n}{m} = \binom{n - 1}{m} + \binom{n - 1}{m - 1}$。

\begin{minted}{cpp}
for (int i = 0; i <= n; i ++) {
    C[i][0] = C[i][i] = 1;
    for (int j = 1; j < i; j ++) {
        C[i][j] = (C[i - 1][j] + C[i - 1][j - 1]) % mod;
    }
}
\end{minted}

\subsubsection{求法 2：预处理阶乘与阶乘逆元}

\begin{quoteblock}
\textbf{适用情况：要求模数为质数，可以承受线性时间的预处理（$n \leq 10^6$）。}

时间复杂度：预处理 $\mathcal{O}(n)$，查询 $\mathcal{O}(1)$。
\end{quoteblock}

将组合数使用阶乘形式表达

\[
\binom{n}{m} = \frac{n!}{m!(n - m)!}
\]
考虑预处理出阶乘与阶乘逆元。

设 $\mathrm{fact}_i = i!$，则有初值 $\mathrm{fact}_0 = 1$，递推式 $\mathrm{fact}_i =\mathrm{fact}_{i - 1} \times i$。

设 $\mathrm{facv}_i = \frac{1}{i!}$，则有初值 $\mathrm{facv}_n = \mathrm{fact}_n^{-1}$，递推式 $\mathrm{facv}_i = \mathrm{facv}_{i + 1} \times(i + 1)$。

于是先正着递推一遍得到 $\mathrm{fact}_{0, \dots, n}$，然后通过一次求逆元得到 $\mathrm{facv}_n$，最后再反着递推一遍得到 $\mathrm{facv}_{n, \dots, 0}$。

\InputCode{cpp}{contents/0x30 数学/code/0x33 组合数（预处理阶乘与阶乘逆元）.cpp}

\subsubsection{求法 3：暴力统计}

\begin{quoteblock}
\textbf{适用范围：要求模数为质数，$m$ 的范围可以承受线性枚举。}

时间复杂度：无预处理，查询 $\mathcal{O}(m)$。
\end{quoteblock}

将组合数表达为

\[
\binom{n}{m} = \frac{n^{\underline{m}}}{m!}
\]
注意到分子分母均只有 $m$ 项，当 $n$ 很大但 $m$ 很小时，可以暴力枚举。

具体实现时，可以先分别求出分子与分母，然后再对分母求一次逆元。

\subsubsection{求法 4：Lucas 定理}

\begin{quoteblock}
\textbf{适用范围：要求模数为质数，模数以内的组合数可以预处理。}

时间复杂度：预处理 $\mathcal{O}(p)$，查询 $\mathcal{O}(\log_{p} n)$。
\end{quoteblock}

由 Lucas 定理，得

\[
\binom{n}{m} \equiv \binom{\lfloor n / p \rfloor}{\lfloor m / p \rfloor}\binom{n \bmod p}{m \bmod p} \pmod p
\]
先预处理出模数以内的组合数。然后不断地使用 Lucas 定理，直到 $n, m$ 递降至模数以内。

\subsubsection{求法 5：exLucas 定理}

\sout{待填坑。}

\subsubsection{求法 6：分段打表}

\begin{quoteblock}
\textbf{适用范围：要求模数为固定质数，代码长度限制不小，可以承受段长 $\mathcal{O}(W)$ 时间的查询。}

\textbf{（正常数据范围 $n \leq 10^9$，配合 Lucas 定理可以做到 $\mathrm{mod} \leq 10^9$ 且 $n \leq 10^{18}$）}

时间复杂度：打表预处理，查询 $\mathcal{O}(W)$。
\end{quoteblock}

仍然将组合数使用阶乘形式表达，每次查询相当于要得到三个阶乘的值。

考虑分段打表。记段长为 $W$（取 $10^5 \sim 10^6$ 相对合适），先通过打表预处理出所有 $W$ 的倍数的阶乘。

每次查询 $n!$ 的时候，就找到最近的 $W$ 的倍数，根据该答案继续向上乘。

\subsubsection{邪修：快速阶乘算法}

\begin{quoteblock}
\textbf{适用范围：要求模数为不固定质数（$n \leq 10^9$）。}

时间复杂度：无预处理，查询 $\mathcal{O}(\sqrt{n} \log n)$。
\end{quoteblock}

直接使用快速阶乘算法！得到三个阶乘的值。

\subsubsection{求法 7：阶乘分解质因数}

\begin{quoteblock}
\textbf{适用范围 1：对模数不做要求，可以承受线性时间的查询（$n \leq 10^6$）。}

时间复杂度：预处理 $\mathcal{O}(n)$，查询 $\mathcal{O}(n)$。

\textbf{适用范围 2：高精度组合数（$n \leq 5000$）。}
\end{quoteblock}

仍然将组合数使用阶乘形式表达，我们考虑 $1 \sim n$ 中的每个质数出现的次数。

先将 $1 \sim n$ 中所有的质数提前筛出来，大约有 $\frac{n}{\ln n}$ 个质数。对于质数 $p$，在阶乘 $n!$ 中 $p$ 的出现次数为

\[
\sum_{i = 1}^{\lfloor \log_p n \rfloor} \left\lfloor \frac{n}{p^i} \right\rfloor
\]
于是将三个阶乘都进行质因数分解，得出每个质数的出现次数，然后将所有质数的贡献相乘即可。

时间复杂度 $\mathcal{O}(n)$（高精度组合数的复杂度另算）。

\subsection{斯特林数}

\subsubsection{第一类斯特林数}

\sout{待填坑。}

\subsubsection{第二类斯特林数}

\sout{待填坑。}

\subsection{卡特兰数}

\Highlight{卡特兰数}：$\mathrm{Cat}_n$，从原点 $(0, 0)$ 出发，每次只能向 $x$ 轴或 $y$ 轴正方向走，且不能越过直线 $y = x$，到达 $(n, n)$ 的方案数。

下标从 $0$ 开始的若干项：

\[
1, 1, 2, 5, 14, 42, 132, 429, 1430, \dots
\]
\Highlight{卡特兰数通项公式}：

\[
\begin{aligned}
\mathrm{Cat}_n & = \binom{2n}{n} - \binom{2n}{n - 1} \\
& = \frac{\binom{2n}{n}}{n + 1}
\end{aligned}
\]
简单推导：越过直线 $y = x$，相当于是触碰直线 $y = x + 1$。

将触碰 $y = x + 1$ 的第一个点以后的路径，关于 $y = x + 1$ 对称，就可以得到一条从 $(0, 0)$ 到 $(n - 1, n + 1)$ 的路径。

\textbf{于是「$(0, 0) \to (n, n)$ 且越过 $y = x$ 的路径」与「$(0, 0) \to (n - 1, n + 1)$ 的路径」是一一对应的}，方案数均为 $\binom{2n}{n - 1}$。

正难则反，不越过直线 $y = x$ 的方案数，即为总方案数 $\binom{2n}{n}$ 减去越过直线的方案数 $\binom{2n}{n - 1}$。

\Highlight{卡特兰数递推公式}：

\[
\begin{gathered}
\mathrm{Cat}_0 = 1 \\
\mathrm{Cat}_n = \sum_{i = 0}^{n - 1}\mathrm{Cat}_i\mathrm{Cat}_{n - i - 1}
\end{gathered}
\]
简单推导：\textbf{枚举第一次触碰直线 $y = x$ 的位置 $(i, i)$}，其中 $i \geq 1$。

由于在 $(i, i)$ 之前不能触碰直线 $y = x$，相当于是 $(1, 0) \to (i, i - 1)$ 不越过直线 $y = x - 1$，于是这部分的方案数为 $\mathrm{Cat}_{i - 1}\mathrm{Cat}_{n - i}$。

\Highlight{卡特兰数递推公式 2}：

\[
\begin{gathered}
\mathrm{Cat}_0 = 1\\
\mathrm{Cat}_n = \frac{4n - 2}{n + 1}\mathrm{Cat}_{n - 1}
\end{gathered}
\]
\Highlight{常见卡特兰数问题}：

\begin{itemize}
\item 从原点 $(0, 0)$ 出发，每次只能向 $x$ 轴或 $y$ 轴正方向走，且不能越过直线 $y = x$，到达 $(n, n)$ 的方案数 $\mathrm{Cat}_n$。
\item 以 $1, 2, \dots, n$ 作为 $n$ 个元素的进栈序列时，出栈序列的方案数 $\mathrm{Cat}_n$。
\item 以 $n$ 对括号组成的合法括号序列的方案数 $\mathrm{Cat}_n$。
\item 以 $n$ 个节点构造无标号、左右儿子区分的二叉树个数 $\mathrm{Cat}_n$。
\item 以 $n$ 个节点构造无标号、儿子区分的有根树个数 $\mathrm{Cat}_{n - 1}$。
\item 以圆周上的 $2n$ 个点为端点，连互不相交的 $n$ 条弦的方案数 $\mathrm{Cat}_n$。
\item 将凸 $n$ 边形用其中 $n - 3$ 条对角线分割为互不重叠的 $n - 2$ 个三角形的方案数 $\mathrm{Cat}_{n - 2}$。
\item 将 $1 \sim 2n$ 中的每个数，不重不漏地填入一个 $2 \times n$ 的矩阵，每个元素大于其左方、上方元素的方案数 $\mathrm{Cat}_n$。
\end{itemize}

\subsection{贝尔数}

\Highlight{贝尔数}：$B_n$，表示将 $n$ 个两两不同的元素，划分成若干个互不区分的非空子集的方案数。

下标从 $0$ 开始的若干项：

\[
\begin{gathered}
1, 1, 2, 5, 15, \\
52, 203, 877, 4140, \\
21147, 115975, 678578, 4213597, \\
27644437, 190899322, 1382958545, \dots
\end{gathered}
\]
\Highlight{贝尔数递推公式}：

\[
\begin{gathered}
B_0 = 1 \\
B_n = \sum_{i = 0}^{n - 1} \binom{n - 1}{i}B_i
\end{gathered}
\]
简单推导：枚举与元素 $n$ 在同一集合的元素个数 $i$，共有 $\binom{n - 1}{i}$ 种选法，剩下的元素继续划分，于是这部分的方案数为 $\binom{n - 1}{i}B_{n - 1 - i}$。

\Highlight{\textbf{贝尔数级别的暴搜}}：当我们需要将 $n$ 个两两不同的元素，划分成若干个互不区分的非空子集时。\textbf{若 $n$ 较小可以直接暴搜，复杂度为 $B_n$ 级别}。

设 $s_1, \dots, s_{t}$ 表示当前时刻已经存在的集合，每次加入一个新元素 $x$ 时有两种转移

\begin{itemize}
\item 选择一个已经存在的集合 $s_i$，将 $x$ 加入 $s_i$。
\item 新开一个集合 $s_{t + 1}$，将 $x$ 加入 $s_{t + 1}$。
\end{itemize}

\Highlight{贝尔数与第二类斯特林数的关系}：

\[
B_n = \sum_{i = 0}^n {n \brace i}
\]
因为第二类斯特林数 ${n \brace i}$ 表示将 $n$ 个两两不同的元素，划分成 $i$ 个互不区分的非空子集的方案数。

\subsection{斐波那契数列}

\Highlight{斐波那契数列}：$\mathrm{Fib}_n$。

\Highlight{斐波那契数列递推公式}：

\[
\begin{gathered}
\mathrm{Fib}_0 = 0, \mathrm{Fib}_1 = 1\\
\mathrm{Fib}_n = \mathrm{Fib}_{n - 1} + \mathrm{Fib}_{n - 2}
\end{gathered}
\]
\Highlight{斐波那契数列递推公式（矩阵形式）}：

\[
\begin{bmatrix}
\mathrm{Fib}_n & \mathrm{Fib}_{n - 1}
\end{bmatrix}
=
\begin{bmatrix}
\mathrm{Fib}_{n - 1} & \mathrm{Fib}_{n - 2}
\end{bmatrix}
\times
\begin{bmatrix}
1 & 1 \\
1 & 0
\end{bmatrix}
\]
即

\[
\begin{bmatrix}
\mathrm{Fib}_n & \mathrm{Fib}_{n - 1}
\end{bmatrix}
=
\begin{bmatrix}
1 & 0
\end{bmatrix}
\times
\begin{bmatrix}
1 & 1 \\
1 & 0
\end{bmatrix}^{n - 1}
\]
该转移矩阵可以求逆，从而实现逆向转移。

\Highlight{斐波那契数列通项公式}：

\[
\mathrm{Fib}_n = \frac{\left( \frac{1 + \sqrt{5}}{2} \right)^n - \left( \frac{1 - \sqrt{5}}{2} \right)^n}{\sqrt{5}}
\]
当 $5$ 是模 $p$ 意义下的\textbf{二次剩余}时（一个常见模数是 $10^9 + 9$），可以直接暴力预先找出 $\sqrt{5}$，然后使用通项公式计算。

否则，考虑使用\textbf{扩展数域}进行计算，定义扩展数域 $\mathbb{F}_p(\sqrt{5})$ 为

\[
\{ a + b\sqrt{5} \mid a, b \in \mathbb{F}_p \}
\]
可以验证，$\mathbb{F}_p(\sqrt{5})$ 关于加、减、乘、除封闭。记 $(a, b)$ 表示 $a + b\sqrt{5}$，则有

\[
\begin{aligned}
(a, b) \pm (c, d) & = (a\pm c, b\pm d) \\
(a, b) \times (c, d) & = (ac + 5bd, ad + bc) \\
(a, b) \div (c, d) & = \left( \frac{ac - 5bd}{c^2 - 5d^2} , \frac{bc - ad}{c^2 - 5d^2} \right)
\end{aligned}
\]
\Highlight{广义斐波那契数列}：将递推公式反过来写，有 $\mathrm{Fib}_n = \mathrm{Fib}_{n + 2} - \mathrm{Fib}_{n + 1}$。

考虑将斐波那契数列扩展到负数域，经过推导可得

\[
\mathrm{Fib}_{-n} = (-1)^{n - 1}\mathrm{Fib}_n
\]
\Highlight{类斐波那契数列}：已知数列 $G_0 = a, G_1 = b$，$G_n = G_{n - 1} + G_{n - 2}$，则有

\[
G_n = \mathrm{Fib}_{n - 1} \times a + \mathrm{Fib}_n \times b
\]
类斐波那契数列相加，仍然是类斐波那契数列。只需将相应的 $a, b$ 相加即可。

\subsubsection{斐波那契数列常见恒等式}

\Highlight{斐波那契数列求和}：

\[
\sum_{i = 1}^n \mathrm{Fib}_i = \mathrm{Fib}_{n + 2} - 1
\]
\Highlight{斐波那契数列平方和}：

\[
\sum_{i = 1}^n \mathrm{Fib}_i^2 = \mathrm{Fib}_n\mathrm{Fib}_{n + 1}
\]
\Highlight{卡西尼性质}：

\[
\mathrm{Fib}_{n - 1}\mathrm{Fib}_{n + 1} - \mathrm{Fib}_n^2 = (-1)^n
\]
简单推导：$\mathrm{Fib}_{n - 1}\mathrm{Fib}_{n + 1} - \mathrm{Fib}_n^2 = \det\begin{pmatrix} \mathrm{Fib}_{n + 1} & \mathrm{Fib}_n \\\mathrm{Fib}_n & \mathrm{Fib}_{n - 1} \end{pmatrix} = \det\begin{pmatrix} 1 & 1 \\1 & 0 \end{pmatrix}^n = (-1)^n$。

\Highlight{附加性质}：

\[
\mathrm{Fib}_{n + m} = \mathrm{Fib}_{n + 1}\mathrm{Fib}_m + \mathrm{Fib}_n\mathrm{Fib}_{m - 1}
\]
简单推导：可以使用 “类斐波那契数列” 的知识解释。

记 $G_0 = \mathrm{Fib}_n, G_1 = \mathrm{Fib}_{n + 1}$，$G_n = G_{n - 1} + G_{n - 2}$，则有 $\mathrm{Fib}_{n + m} = G_m$。

\Highlight{斐波那契数列相邻项互质性质}：

\[
\gcd(\mathrm{Fib}_n, \mathrm{Fib}_{n + 1}) = 1
\]
\Highlight{斐波那契数列 gcd 性质}：

\[
\gcd(\mathrm{Fib}_n, \mathrm{Fib}_m) = \mathrm{Fib}_{\gcd(n, m)}
\]
简单推导：不妨设 $n < m$，则有

\[
\begin{aligned}
\gcd(\mathrm{Fib}_n, \mathrm{Fib}_m) & = \gcd(\mathrm{Fib}_n, \mathrm{Fib}_{n + (m - n)}) \\
& = \gcd(\mathrm{Fib}_n, \mathrm{Fib}_{n + 1}\mathrm{Fib}_{m - n} + \mathrm{Fib}_n\mathrm{Fib}_{m - n - 1}) \\
& = \gcd(\mathrm{Fib}_n, \mathrm{Fib}_{n + 1}\mathrm{Fib}_{m - n})
\end{aligned}
\]
由上文提到的 $\mathrm{Fib}_{n + 1}$ 与 $\mathrm{Fib}_n$ 互质。则有 $\gcd(\mathrm{Fib}_n, \mathrm{Fib}_m) = \gcd(\mathrm{Fib}_n, \mathrm{Fib}_{m - n})$。该式子与辗转相除法相似，得证。

\Highlight{斐波那契数列整除性质}：

\[
\mathrm{Fib}_n \mid \mathrm{Fib}_m \iff n \leq 2 \lor n \mid m
\]

\subsubsection{斐波那契数列模意义下的循环节}

在模意义下，二元组 $(\mathrm{Fib}_n \bmod p, \mathrm{Fib}_{n + 1} \bmod p)$ 仅有 $p^2$ 种取值。\textbf{因此，当 $n$ 足够大时，必定会出现循环节}。

根据递推式，有 $\mathrm{Fib}_n = \mathrm{Fib}_{n + 2} - \mathrm{Fib}_{n + 1}$，\textbf{也就是说 $\mathrm{Fib}_{n + 1}, \mathrm{Fib}_{n + 2}$ 的前面必为 $\mathrm{Fib}_n$}。

根据这两个事实我们可以得知，\textbf{斐波那契数列在模意义下必有循环节，并且一定是纯循环的（循环节的起点为 $\mathrm{Fib}_0, \mathrm{Fib}_1$）}。

\Highlight{最小循环节（皮萨诺周期）}：定义斐波那契数列模 $p$ 意义下的最小循环节 $\pi(p)$ 为最小的正整数 $m$，使得 $\mathrm{Fib}_m \equiv 0 \pmod p$ 且 $\mathrm{Fib}_{m + 1} \equiv 1 \pmod p$。

注意：周期一定是最小循环节的倍数，即

\[
\begin{cases} \mathrm{Fib}_k \equiv 0 \pmod p \\ \mathrm{Fib}_{k + 1} \equiv 1 \pmod p \end{cases} \iff \pi(p) \mid k
\]
很多时候我们只需求出周期即可。

\Highlight{最小循环节的上界}：\textbf{斐波那契数列 $\operatorname{mod} m$ 的最小正周期不超过 $6m$，即 $\pi(m) \leq 6m$}。

\Highlight{求斐波那契数列循环节（结论法）}：

\begin{itemize}
\item 质数情景：
  \begin{itemize}
  \item $\pi(2) = 3$，$\pi(5) = 20$。
  \item 当 $p$ 为奇质数且 $p \equiv \pm 1 \pmod 5$ 时，$\pi(p) \mid p - 1$。
  \item 当 $p$ 为奇质数且 $p \equiv \pm 2 \pmod 5$ 时，$\pi(p) \mid 2p + 2$。
  \end{itemize}
\item 一般情景：
  \begin{itemize}
  \item 对于质数 $p$ 以及任意正整数 $k$，有 $\pi(p^k) \mid p^{k - 1}\pi(p)$。
  \item 对于 $m = \prod p_i^{c_i}$，有 $\pi(m) = \mathrm{lcm}\{ \pi(p_i^{c_i}) \}$，同时有 $\pi(m) \mid \mathrm{lcm}\{ p_i^{c_i - 1}\pi(p_i) \}$。
  \end{itemize}
\end{itemize}

在一般应用中，通常只需要求出一个循环节即可。

对于质数 $p$，记

\[
g(p) = \begin{cases} 3 & p = 2 \\
20 & p = 5 \\
p - 1 & p \equiv \pm 1 \pmod 5\\
2p + 2 & p \equiv \pm 2 \pmod 5
\end{cases}
\]
然后将 $m$ 质因数分解成 $\prod p_i^{c_i}$，使用 $\mathrm{lcm}\{ p_i^{c_i - 1}g(p_i) \}$ 作为循环节即可（注意这并不是最小循环节，这个值在多数情况下已经够用了）。

\Highlight{求斐波那契数列循环节（矩阵 BSGS）}：

\begin{quoteblock}
\textbf{适用场景：最小循环节的上界确定，转移矩阵可逆。}
\end{quoteblock}

考虑斐波那契数列的矩阵表示

\[
\begin{bmatrix}
\mathrm{Fib}_n & \mathrm{Fib}_{n - 1}
\end{bmatrix}
=
\begin{bmatrix}
1 & 0
\end{bmatrix}
\times
\begin{bmatrix}
1 & 1 \\
1 & 0
\end{bmatrix}^{n - 1}
\]
记转移矩阵 $A = \begin{bmatrix} 1 & 1 \\ 1 & 0 \end{bmatrix}$，那么循环节即为满足 $A^m \equiv I \pmod p$ 的正整数 $m$。

在转移矩阵可逆的情况下，套用 BSGS 的流程：取 $S$ 表示上界的根号，设 $m = iS - j$（其中 $i \geq 1$ 且 $0 \leq j < S$），则有

\[
\begin{gathered}
A^{iS} \times A^{-j} \equiv I \pmod p \\
A^{iS} \equiv A^j \pmod p
\end{gathered}
\]
预处理 $0 \leq j < S$ 的 $A^j$ 存入哈希表，每次在哈希表中查询 $A^{iS}$ 即可。

时间复杂度 $\mathcal{O}(\sqrt{p})$。

\Highlight{求斐波那契数列循环节（生日悖论随机化）}：

若存在两个数 $0 \leq i, j \leq 6p$，满足 $(\mathrm{Fib}_i, \mathrm{Fib}_{i + 1}) \equiv (\mathrm{Fib}_j, \mathrm{Fib}_{j + 1}) \pmod p$，则 $|i - j|$ 为一个循环节（注意这并不是最小循环节）。

每次新随机一个数 $x$，将二元组 $(\mathrm{Fib}_x, \mathrm{Fib}_{x + 1})$ 存入哈希表中，并查询哈希表中有没有与之相同的 $(\mathrm{Fib}_y, \mathrm{Fib}_{y + 1})$。

发现这与生日悖论的形式很像，于是时间复杂度正确。

\begin{quoteblock}
生日悖论：在 $[1, n]$ 中随机选择正整数，选出两个相同数的期望轮数为 $\mathcal{O}(\sqrt{n})$ 级别（近似 $\sqrt{\frac{\pi n}{2}}$）。
\end{quoteblock}

实际取数过程可能需要在 $[0, 12p]$ 中随机以保证期望次数。求 $\mathrm{Fib}_{i}$ 需要用到 “光速幂” 的技巧。

\subsection{普通生成函数}

\Highlight{普通生成函数（OGF）}：

\[
F(x) = \sum_{n} a_nx^n
\]
\Highlight{OGF 乘法}：

\[
F(x)G(x) = \sum_{n} \left( \sum_{i = 0}^n a_ib_{n - i} \right)x^n
\]
OGF 常用于解决\textbf{无标号的计数问题}（例如：优化计数类背包）。

\Highlight{常见序列 OGF 封闭形式}：

\begin{itemize}
\item $\{1, 1, 1, \dots \} \xrightarrow{\mathrm{OGF}} \frac{1}{1 - x}$
\item $\{1, p, p^2, \dots \} \xrightarrow{\mathrm{OGF}} \frac{1}{1 - px}$
\item $\{ \binom{n}{0}, \binom{n}{1}, \binom{n}{2}, \dots \} \xrightarrow{\mathrm{OGF}} (1 + x)^n$
\item $\{ \binom{n}{0}, \binom{n + 1}{1}, \binom{n + 2}{2}, \dots \} \xrightarrow{\mathrm{OGF}} \frac{1}{(1 - x)^{n + 1}}$
\item $\{ 0, 1, \frac{1}{2}, \frac{1}{3}, \frac{1}{4}, \dots \} \xrightarrow{\mathrm{OGF}} \ln \frac{1}{1 - x} = -\ln(1 - x)$
\item $\{ 1, 1, \frac{1}{2!}, \frac{1}{3!}, \frac{1}{4!}, \dots \} \xrightarrow{\mathrm{OGF}} e^x$
\end{itemize}

\subsection{指数生成函数}

\Highlight{指数生成函数（EGF）}：

\[
F(x) = \sum_{n} a_n \frac{x^n}{n!}
\]
\Highlight{EGF 乘法}：

\[
F(x)G(x) = \sum_{n} \left( \sum_{i = 0}^n \binom{n}{i}a_ib_{n - i} \right)\frac{x^n}{n!}
\]
EGF 常用于解决\textbf{有标号的计数问题}。

\Highlight{常见序列 EGF 封闭形式}：

\begin{itemize}
\item $\{1, 1, 1, \dots\} \xrightarrow{\mathrm{EGF}} e^x$
\item $\{1, p, p^2, \dots\} \xrightarrow{\mathrm{EGF}} e^{px}$
\item $\{1, n, n^{\underline{2}}, n^{\underline{3}}, n^{\underline{4}} \dots \} \xrightarrow{\mathrm{EGF}} (1 + x)^n$
\end{itemize}

\Highlight{EGF 中多项式 exp 的组合意义}：

回顾一下 exp 的式子（在这里我们默认 $[x^0]f(x) = 0$）

\[
\exp f(x) = \sum\limits_{i = 0}^{\infty} \frac{f^i(x)}{i!}
\]
如果 $f(x)$ 描述的是单个组件的 EGF，那么 $f^k(x)$ 得到的就是由 $k$ 个组件构成的有序序列的 EGF，但是如果我们将这 $k$ 个组件看成一个集合，就要除以 $k!$ 来消除顺序的影响。

\textbf{所以 $\frac{f^k(x)}{k!}$ 的含义是由 $k$ 个组件构成的无序集合的 EGF}。因此，\textbf{$\exp f(x)$ 的含义是由若干个组件构成的无序集合的 EGF}。

\subsection{符号化方法}

\sout{待填坑。}

\subsection{图论计数}

\subsubsection{有标号无向连通图计数}

\paragraph{算法一：分治 NTT}

设 $f_i$ 表示点数为 $i$ 的\textbf{有标号简单无向连通图}个数，设 $g_i$ 表示点数为 $i$ 的\textbf{有标号简单无向图}个数，显然 $g_i = 2^{\binom{i}{2}}$。

考虑容斥。枚举点 $1$ 所在的连通块大小，有递推式

\[
\begin{gathered}
f_1 = 1 \\
f_i = g_i - \sum_{j = 1}^{i - 1}\binom{i - 1}{j - 1}f_jg_{i - j}
\end{gathered}
\]
化简式子，得

\[
\frac{f_i}{(i - 1)!} = \frac{g_i}{(i - 1)!} - \sum_{j = 1}^{i - 1}\frac{f_j}{(j - 1)!} \frac{g_{i - j}}{(i - j)!}
\]
使用分治 NTT，时间复杂度 $\mathcal{O}(n \log^2 n)$。

\begin{quoteblock}
另一种理解方式：设 $F_i = \frac{f_i}{(i - 1)!}, G_i = \frac{g_i}{(i - 1)!}, H_i = \frac{g_i}{i!}$，可以将该递推式看成 $G = F \times H$，于是 $F = G \times H^{-1}$。

使用多项式求逆，时间复杂度 $\mathcal{O}(n \log n)$。
\end{quoteblock}

\paragraph{算法二：多项式 ln}

设 $F(x)$ 是 $f_i$ 的 EGF，$G(x)$ 是 $g_i$ 的 EGF。由 EGF 中多项式 exp 的组合意义，有 $G(x) = e^{F(x)}$。所以

\[
F(x) = \ln G(x)
\]
使用多项式 ln，时间复杂度 $\mathcal{O}(n \log n)$。

\InputCode{cpp}{contents/0x30 数学/code/0x33 有标号无向连通图计数.cpp}

\subsubsection{有标号无向欧拉图计数}

欧拉图的条件：所有节点的度数均为偶数，并且连通。

设 $f_i$ 表示点数为 $i$ 的\textbf{有标号简单无向欧拉图}个数，设 $g_i$ 表示点数为 $i$ 的\textbf{有标号度数均为偶数的简单图}个数。

考虑求 $g_i$，先让前 $i - 1$ 个点任意连边，最后所有奇点的度数都通过与 $i$ 连边调节成偶数。于是 $g_i = 2^{\binom{i - 1}{2}}$。

剩余部分同 “有标号无向连通图计数”，时间复杂度 $\mathcal{O}(n \log n)$。

\InputCode{cpp}{contents/0x30 数学/code/0x33 有标号无向欧拉图计数.cpp}

\subsubsection{有标号无向二分图计数}

设 $f_i$ 表示点数为 $i$ 的\textbf{有标号简单无向二分图}个数，设 $g_i$ 表示点数为 $i$ 的\textbf{对节点进行 $2$ 染色，满足同色节点不存在边的图（二分染色图）}个数。

考虑求 $g_i$，枚举黑色节点个数 $j$，则有

\[
\begin{gathered}
g_i = \sum_{j = 0}^i \binom{i}{j} 2^{j(i - j)} \\
g_i = \sum_{j = 0}^i \binom{i}{j} 2^{\binom{i}{2} - \binom{j}{2} - \binom{i - j}{2}} \\
\frac{g_i}{i! \cdot 2^{\binom{i}{2}}} = \sum_{j = 0}^i \frac{1}{j! \cdot 2^{\binom{j}{2}}} \cdot \frac{1}{(i - j)! \cdot 2^{\binom{i - j}{2}}}
\end{gathered}
\]

\begin{quoteblock}
这里用到了 $ij = \binom{i + j}{2} - \binom{i}{2} - \binom{j}{2}$。
\end{quoteblock}

可以使用卷积 $\mathcal{O}(n \log n)$ 求出 $g_i$。

对于二分染色图的每一个连通块，确定了最小点的颜色就确定了整个连通块的颜色，而最小点有黑白两种选法。

则二分染色图可以由两个组件构成，第一个组件是最小点颜色为黑色的连通块集合，第二个组件是最小点颜色为白色的连通块集合。而两个组件内部均是二分图。

于是设 $F(x)$ 是 $f_i$ 的 EGF，$G(x)$ 是 $g_i$ 的 EGF，则有

\[
\begin{gathered}
F(x) \times F(x) = G(x) \\
F(x) = \sqrt{G(x)}
\end{gathered}
\]
使用多项式开方，时间复杂度 $\mathcal{O}(n \log n)$。

\InputCode{cpp}{contents/0x30 数学/code/0x33 有标号无向二分图计数.cpp}

\subsubsection{有标号 DAG 计数}

设 $f_i$ 表示点数为 $i$ 的\textbf{有标号 DAG（不保证弱连通）}个数。

将 DAG 视作一个分层图，考虑容斥，\textbf{钦定}拓扑序第一层（零入度点）的个数。

\[
f_i = \sum_{j = 1}^{i}(-1)^{j + 1} \binom{i}{j} 2^{j(i - j)}f_{i - j}
\]
至此，可以做到 $\mathcal{O}(n^2)$。

进一步

\[
\frac{f_i}{i! \cdot 2^\binom{i}{2}} = \sum_{j = 1}^i \frac{(-1)^{j + 1}}{j! \cdot 2^\binom{j}{2}} \cdot \frac{f_{i - j}}{(i - j)! \cdot 2^\binom{i - j}{2}}
\]
设 $F_i = \frac{f_i}{i! \cdot 2^\binom{i}{2}}, G_j = \frac{(-1)^{j + 1}}{j! \cdot 2^\binom{j}{2}}$，则有

\[
\begin{gathered}
F = G\times F + 1 \\
F = \frac{1}{1 - G}
\end{gathered}
\]
使用多项式求逆，时间复杂度 $\mathcal{O}(n \log n)$。

\InputCode{cpp}{contents/0x30 数学/code/0x33 有标号 DAG 计数.cpp}

\subsubsection{有标号弱连通 DAG 计数}

设 $f_i$ 表示点数为 $i$ 的\textbf{有标号 DAG（不保证弱连通）}个数，设 $h_i$ 表示点数为 $i$ 的\textbf{有标号弱连通 DAG} 个数。$f_i$ 的求解详见 “有标号 DAG 计数”。

设 $F(x)$ 是 $f_i$ 的 EGF，$H(x)$ 是 $h_i$ 的 EGF。由 EGF 中多项式 exp 的组合意义，有 $F(x) = e^{H(x)}$。所以

\[
H(x) = \ln F(x)
\]
使用多项式 ln，时间复杂度 $\mathcal{O}(n \log n)$。

\InputCode{cpp}{contents/0x30 数学/code/0x33 有标号弱连通 DAG 计数.cpp}

\subsubsection{有标号强连通图计数}

设 $f_i$ 表示点数为 $i$ 的\textbf{有标号强连通图}个数，设 $b_{i, j}$ 表示 \textbf{$i$ 个点形成 $j$ 个独立的强连通图}的个数，设 $h_i$ 表示点数为 $i$ 的\textbf{任意有向图}个数。显然 $h_i = 2^{i(i - 1)}$。

则有

\[
h_i = \sum_{j = 1}^i \binom{i}{j} h_{i - j}\times 2^{j(i - j)} \times \sum_{k = 1}^{j}(-1)^{k + 1}b_{j, k}
\]
对于任意点数为 $i$ 的有向图，将每个强连通分量缩成一个点，会得到一张 DAG。这其中必定有出度为 $0$ 的点。

\textbf{钦定 $k$ 个强连通分量出度为 $0$，这 $k$ 个强连通分量里包含了 $j$ 个点}。剩下的 $i - j$ 个点随意连接。

设 $g_j = \sum_{k = 1}^j (-1)^{k + 1}b_{j, k}$，则有

\[
\begin{gathered}
g_i = h_i - \sum_{j = 1}^{i - 1} \binom{i}{j} 2^{j(i - j)} h_{i - j} g_j \\
\frac{h_i}{i! \cdot 2^{\binom{i}{2}}} = \sum_{j = 1}^i \frac{g_j}{j! \cdot 2^{\binom{j}{2}}} \cdot \frac{h_{i - j}}{(i - j)! \cdot 2^{\binom{i - j}{2}}}
\end{gathered}
\]
设 $H_i = \frac{h_i}{i! \cdot 2^{\binom{i}{2}}}, G_i = \frac{g_i}{i! \cdot 2^{\binom{i}{2}}}$，则有

\[
\begin{gathered}
H = G \times H + 1 \\
G = 1 - \frac{1}{H}
\end{gathered}
\]
至此，可以使用多项式求逆求出 $G$。

再根据 $g$ 的定义，考虑 $g$ 的转移，有

\[
\begin{gathered}
g_i = f_i - \sum_{j = 1}^{i - 1} \binom{i - 1}{j - 1}f_jg_{i - j} \\
f_i = g_i + \sum_{j = 1}^{i - 1} \binom{i - 1}{j - 1}f_jg_{i - j} \\
\frac{f_i}{(i - 1)!} = \frac{g_i}{(i - 1)!} + \sum_{j = 1}^{i - 1}\frac{f_j}{(j - 1)!}\cdot \frac{g_{i - j}}{(i - j)!}
\end{gathered}
\]
设 $F_i = \frac{f_i}{(i - 1)!}, A_i = \frac{g_i}{(i - 1)!}, B_i = \frac{g_i}{i!}$，则有

\[
\begin{gathered}
F = A + F \times B \\
F = \frac{A}{1 - B}
\end{gathered}
\]
\InputCode{cpp}{contents/0x30 数学/code/0x33 有标号强连通图计数.cpp}

\subsection{Misc}

\subsubsection{Raney 引理}

\Highlight{Raney 引理}：若整数序列 $x_1, \dots, x_m$ 满足 $\sum_{i = 1}^m x_i = 1$，则它的 $m$ 个循环移位中恰好有 $1$ 个满足前缀和均为正数。

唯一合法的循环移位起点，即为最靠后的 “前缀和最小点”。

\Highlight{广义 Raney 引理}：若整数序列 $x_1, \dots, x_m$ 满足 $x_i \leq 1$ 且 $\sum_{i = 1}^m x_i = n$，则它的 $m$ 个循环移位中恰好有 $n$ 个满足前缀和均为正数。
